\section{Introduction}\label{sec:introduction}
Computing $\rtwo$ is an elementary problem with a useful lesson about numerical
algorithms: an update rule, a convergence proof, and an implementation are three
different objects. The polynomial $f(x)=x^2-2$ has a unique positive zero, and
both Newton's method and bisection locate it. Their error reduction mechanisms
are very different. Newton uses a derivative to make rapid local progress;
bisection maintains an interval guaranteed to contain a root.

The purpose of this note is to make those mechanisms explicit without relying
on a general convergence theorem. We work in exact real arithmetic for the
analysis, then use binary64 arithmetic for a small illustration. Standard
background on root finding can be found in \cite{burden2015,stoer2002}.
Theorem~\ref{thm:convergence} establishes the main claim, and
Section~\ref{sec:error} explains the quadratic rate.

The example also serves as a compact research-project maintenance exercise.
Its notation, cross-references, bibliography, source code, and derived figures
must stay synchronized. A clean compilation checks document construction;
it does not certify any mathematical assertion.
